\subsection{Quantifiers}
\newcommand{\norforall}[2]{{\forall}#1#2}
\newcommand{\norin}[2]{{\in}#1#2}
\begin{notation}
    \begin{align*}
        \forall a \in b: P &\coloneqq \forall a: (a \in b \Rightarrow P) \\
        \forall a, b \in c: P &\coloneqq \forall a \in c: \forall b \in c: P \\
        \forall a \in b; c \in d: P &\coloneqq \forall a \in b: \forall c \in d: P
    \end{align*}
    \[[\norforall{a}{\norimplies{\norin{a}{b}}{P}}]\]
\end{notation}

\newcommand{\forallI}{{\forall}I}
\newcommand{\forallE}{{\forall}E}
\begin{foundation}
    % Seriously thinking of adopting a quick def for \langle x \rangle - rustlers joy. DONE!
    \[\inferrule{P\bound{a}}{\forall a: P\bound{a}} \quad (\forallI)\]
    \[\inferrule{\forall a: P\bound{a}}{P\bound{b}} \quad (\forallE)\]
\end{foundation}

% Aa:P<=>Q > Aa:P <=> Aa:Q
\begin{align*}
    \premise{\forall a: (P\bound{a} \lequiv Q\bound{a})}\\
    \assumption[1]{\forall a: P\bound{a}}\\
    \byrule[1]{P\bound{b}}{\forallE}\\
    \byrule[1]{P\bound{b} \lequiv Q\bound{b}}{\forallE}\\
    \byrule[1]{Q\bound{b}}{\lequivE}\\
    \byrule[1]{\forall a: Q\bound{a}}{\forallI}\\
    \byrule{(\forall a: P\bound{a}) \implies (\forall a: Q\bound{a})}{\impliesI}\\
    \assumption[1]{\forall a: Q\bound{a}}\\
    \byrule[1]{Q\bound{b}}{\forallE}\\
    \byrule[1]{Q\bound{b} \lequiv P\bound{b}}{\forallE}\\
    \byrule[1]{P\bound{b}}{\lequivE}\\
    \byrule[1]{\forall a: P\bound{a}}{\forallI}\\
    \byrule{(\forall a: Q\bound{a}) \implies (\forall a: P\bound{a})}{\impliesI}\\
    \byrule{(\forall a: P\bound{a}) \lequiv (\forall a: Q\bound{a})}{\lequivI}
\end{align*}
\newcommand{\forallC}{{\forall}\text{Cong}}
\newcommand{\forallD}{{\forall}\text{Cong}} % Deprecated
\[\inferrule{\forall a: (P\bound{a} \lequiv Q\bound{a})}{(\forall a: P\bound{a}) \lequiv (\forall a: Q\bound{a})} \quad (\forallC)\]


%\begin{notation}
%    \begin{align*}
%        \exists a: P &\coloneqq \neg(\forall a: \neg P) \\ % This should be the earliest point that we define existence.
%        \exists a \in b: P &\coloneqq \exists a: (a \in b \wedge P) \\
%        \exists a, b \in c: P &\coloneqq \exists a \in c: \exists b \in c: P \\
%        \exists a \in b; c \in d: P &\coloneqq \exists a \in b: \exists c \in d: P
%    \end{align*}
%\end{notation}
%% All of this again to introduce punctuation intuitively.
%\begin{foundation}
%    %\[\inferrule{P\bound{a}}{\exists b: P\bound{b}} \quad (\existsI)\]
%\end{foundation}
%\[\inferrule{\exists a: P\bound{a} \\ [P\bound{b}] \cdots Q}{Q} \quad (\existsE)\]
% Any notion of bound and free variables above is implicit - no further explanation of these concepts is provided.
